% C25 paradoxo COVID
\begin{table}[htbp]\centering
\caption{Paradoxo COVID: comparac\~ao direta entre INEP e PNADC v5 nos anos 2018--2021. Durante 2020 e 2021, o INEP registra promoc\~ao mais alta e repet\^encia mais baixa (compat\'ivel com aprovac\~ao autom\'atica), enquanto a PNADC registra o oposto.}
\label{tab:paradoxo_covid}
\begin{threeparttable}\small
\begin{tabular}{lrrrrrrrr}\toprule
& \multicolumn{2}{c}{Promo\c{c}\~ao (\%)} & \multicolumn{2}{c}{Repet\^encia (\%)} & \multicolumn{2}{c}{Evas\~ao (\%)} & \multicolumn{2}{c}{$\Delta$ prom.\ (pp)} \\
\cmidrule(lr){2-3} \cmidrule(lr){4-5} \cmidrule(lr){6-7} \cmidrule(lr){8-9}
Ano $t$ & INEP & PNADC & INEP & PNADC & INEP & PNADC & 2020/19 & 2021/19 \\
\midrule
\multicolumn{9}{l}{\textit{EF iniciais}} \\
\midrule
2018 & 92.6 & 86.3 & 6.1 & 11.7 & 1.2 & 0.4 & &  \\
2019 & 93.4 & 86.1 & 5.2 & 11.9 & 1.3 & 0.4 & &  \\
2020 & 96.2 & 76.9 & 2.2 & 21.9 & 1.5 & 0.4 & I:+2.8 P:-9.2 & \\
2021 & 95.5 & 82.9 & 3.1 & 16.0 & 1.3 & 0.2 & & I:+2.1 P:-3.2 \\
\\[0.3em]
\multicolumn{9}{l}{\textit{EF finais}} \\
\midrule
2018 & 85.0 & 81.0 & 8.8 & 13.6 & 3.9 & 2.0 & &  \\
2019 & 87.2 & 81.4 & 8.0 & 13.8 & 3.0 & 1.3 & &  \\
2020 & 93.9 & 73.1 & 2.3 & 24.7 & 2.9 & 0.8 & I:+6.7 P:-8.3 & \\
2021 & 91.2 & 80.4 & 3.3 & 15.9 & 4.0 & 2.0 & & I:+4.0 P:-1.0 \\
\\[0.3em]
\multicolumn{9}{l}{\textit{EM}} \\
\midrule
2018 & 79.5 & 75.7 & 8.2 & 13.6 & 9.7 & 7.4 & &  \\
2019 & 82.7 & 75.8 & 8.3 & 15.0 & 6.9 & 5.9 & &  \\
2020 & 89.0 & 63.5 & 3.9 & 29.9 & 5.9 & 3.8 & I:+6.3 P:-12.3 & \\
2021 & 85.3 & 73.4 & 4.2 & 17.9 & 8.7 & 6.8 & & I:+2.6 P:-2.4 \\
\bottomrule\end{tabular}
\begin{tablenotes}\footnotesize
\item \textit{Fontes:} INEP Indicadores de Trajet\'oria (Brasil) e painel longitudinal PNADC v5.
\item \textit{Notas:} ``$\Delta$ prom.\ 2020/19'' \'e a varia\c{c}\~ao da promoc\~ao em 2020 em relac\~ao a 2019, pelas duas fontes (I = INEP, P = PNADC). Sinais opostos entre INEP e PNADC indicam que as duas fontes capturam dimens\~oes distintas do mesmo fen\^omeno: registro administrativo de matr\'icula vs.\ auto-relato familiar sobre a situac\~ao escolar.
\end{tablenotes}\end{threeparttable}\end{table}
